\documentclass[11pt]{article}
\usepackage[a4paper,margin=24mm]{geometry}
\usepackage{amsmath,amssymb,amsthm,booktabs}
\newtheorem{lemma}{Lemma}\newtheorem{theorem}{Theorem}
\title{Five explicit Higman--Sims doublings\\and three source-link qualifications}
\author{L091 finite construction component}\date{}
\begin{document}\maketitle
We use the archived supplement dated 6 October 2003 to D.~Leemans
\cite{HS}, with its accepted coordinates in the 86-class literal coset
catalogue. All labels below are supplement labels. The main paper calls
the construction of \cite[Theorem 4.1]{LPP} doubling. We use precisely the
two-point-shadow operation in the acquired theorem adapter: the prior
primary reading included that theorem and its construction, but did not
obtain the complete article. No different convention is silently imposed.
The elementary argument needed here is given in full.

\begin{theorem}
The following five printed links have explicit incidence isomorphisms
from the indicated doubled parents to the target geometries:
\[
3.2\longrightarrow3.3,\quad3.4\longrightarrow3.5,\quad
3.8\longrightarrow3.9,\quad4.8\longrightarrow4.10,\quad
4.14\longrightarrow4.19.
\]
The other three printed links, $3.12\to3.15$, $3.13\to3.16$ and
$3.20\to3.21$, are not applications of this operation to the stated
parents. The first two fail its two-point hypothesis at the only type
choice compatible with the target sizes; the third fails the necessary
type-size condition for every choice.
\end{theorem}
These are ancillary construction qualifications, relative to the explicit
operation below. They neither refute the principal classification nor
claim that the three target geometries do not exist. Their existence and
coordinates remain acquired. No repair using an unspecified alternative
doubling convention or an unverified replacement parent is asserted.

\section{The construction and exact finite test}
\begin{lemma}
Suppose a flag-transitive geometry has types $X,E$ such that the
$X$-shadows of its $E$-objects are distinct two-element sets. Suppose, for
any other object $z$,
\[
 e*z\quad\Longleftrightarrow\quad \sigma_X(e)\subseteq\sigma_X(z). \tag{1}
\]
Replace $X,E$ by two copies $X^+,X^-$ of $X$. Join $x^+$ to $y^-$
exactly when $\{x,y\}$ is an $E$-shadow; join either copy of $x$ to
$z$ exactly when $x*z$ originally. Leave other incidences unchanged.
This is a geometry with the same chamber action as the original one.
\end{lemma}
\begin{proof}
An old chamber $(x,e,F)$ determines the new chamber $(x^+,y^-,F)$,
where $y$ is the other endpoint of $e$. Conversely a new flag using both
copies determines its unique edge, and (1) makes that edge incident with
all its unchanged objects. It therefore extends from an old flag to an
old chamber and then to a new chamber. A flag using only one copy extends
from the corresponding old point flag; if neither copy occurs, extend its
unchanged flag first. This proves chamber extension in every case and
the asserted equivariant chamber bijection. Chamber transitivity then
implies flag transitivity. The argument also works in rank two; it does
not assert extra rank-two connectivity hypotheses from a theorem stated
only in higher rank.
\end{proof}

Write $G=\mathrm{HS}$. Objects are right cosets; incidence means nonempty
intersection. We also denote the base point and edge stabilizers by
$X,E$. The number of points incident with $E$ is $[E:E\cap X]$.
When this is two, choose a recorded $h\in E\setminus X$; the endpoints
are $X$ and $Xh$. Their pointwise stabilizer is $X\cap X^h$, where
$X^h=h^{-1}Xh$. Hence distinct edge shadows are equivalent to
\[
                   2|X\cap X^h|=|E|.                    \tag{2}
\]
Indeed, $E$ is contained in the pair's setwise stabilizer and interchanges
its endpoints. That setwise stabilizer has twice the order of the
pointwise stabilizer, so (2) gives equality. The resulting map from
$E$-cosets to pairs is therefore injective.

For any remaining base stabilizer $Z$, take right transversals
$T_X$ of $X\cap Z$ in $X$ and $T_E$ of $E\cap Z$ in $E$.
The following finite conditions are equivalent to (1):
\begin{align}
 h a^{-1}\in XZ&\quad\Longleftrightarrow\quad a^{-1}\in EZ
                                      &&(a\in T_X),\tag{3}\\
 a^{-1}\in XZ\ \text{ and }\ h a^{-1}\in XZ
                                      &&&(a\in T_E).\tag{4}
\end{align}
Here $XZ$ and $EZ$ are identity double cosets, not asserted subgroups.
The objects $Za$, $a\in T_X$, are exactly the $Z$-objects incident
with $X$. The second endpoint meets $Za$ precisely when
$ha^{-1}\in XZ$, while $E$ meets it precisely when $a^{-1}\in EZ$.
This proves sufficiency of a pair of endpoints in (3). Equation (4)
checks necessity for every $Z$-neighbor of $E$, including those not
assumed to meet $X$. Transitivity transports these tests to all edges.
No enumeration of the $44\,352\,000$ ambient elements is needed.

\section{Five literal target maps}
The new geometry has base stabilizers
\[
                     X,\ X^h,\ (Z_i)_{i\ne X,E}.        \tag{5}
\]
In the minus copy the base point is $Xh$, so the right coset $X^h g$
maps to $Xhg$. The chamber argument makes every type-pair incidence a
single $G$-orbit, and thus (5) is its full coset description.
Each certificate supplies an actual $n$ in the acquired normalizer
$N$ of $G$ and a type permutation, and checks that conjugating every
subgroup in (5) by $n$ gives the target family in that order.
Consequently $Hg\mapsto H^n g^n$ is a well-defined typewise bijection
preserving incidence and nonincidence. This is an exact incidence map,
not an inference from a diagram or subgroup order.

The following table uses internal catalogue type indices solely to make
the literal certificates easy to locate. The source-coordinate ledger
also supplies the corresponding native node indices.
\begin{center}\small
\begin{tabular}{ccccc}\toprule
Parent$\to$target & Internal parent$\to$target & $(X,E)$ & (3) tests & (4) tests\\\midrule
$3.2\to3.3$ & $35\to34$ & $(3,1)$ &8&2\\
$3.4\to3.5$ & $32\to33$ & $(3,2)$ &15&3\\
$3.8\to3.9$ & $30\to28$ & $(3,1)$ &8&2\\
$4.8\to4.10$ & $70\to69$ & $(4,3)$ &2310&48\\
$4.14\to4.19$ & $54\to57$ & $(2,3)$ &2800&55\\\bottomrule
\end{tabular}
\end{center}
The two rank-four rows test both unchanged types. The accepted target
geometries transport all their residue properties through these maps.
The present component does not infer full automorphism or correlation
groups merely from the induced specified $G$-action.

\section{Three exact obstructions and evidence scope}
For the links $3.12\to3.15$ and $3.13\to3.16$, the parent and target
size lists, in the same order, are
\[
\begin{array}{ccl}
(352,30800,1100)&\longrightarrow&(352,1100,1100),\\
(352,30800,5600)&\longrightarrow&(352,5600,5600).
\end{array}
\]
Thus the only possible choices duplicate the 1100- or 5600-object type
and remove the 30800-object type. In each case the edge stabilizer has
order $1440$, and its intersection with the proposed point stabilizer has
order $144$. Each edge shadow therefore has ten points, not two. These
are literal subgroup intersections, also acquired in the numerical
metadata. In particular a type permutation cannot fix the obstruction.

For $3.20\to3.21$, the parent and target size lists are
\[
(15400,346500,15400)\quad\hbox{and}\quad(30800,15400,15400).
\] The operation replaces two sizes by two copies
of one existing size and retains the third. It cannot create the new
size $30800$. The native statements are on supplement pages6,7 and8;
the eight links and their source-node maps are explicitly retained.

The checker reads the exact previously accepted group, complete family
and source-coordinate inputs. GAP's finite subgroup, transversal and
permutation computations remain explicit computational premises. It
verifies all five saved subgroup maps directly, all local tests (2)--(4),
and all three obstructions. Thirty damaged certificates are rejected.
An independent literal-coset control over $S_3,D_8,A_4$ compares the
criterion with all actual pair shadows and all relevant edge--object
incidences, including both successful and failed endpoint conditions.
The failed first probe is preserved: it returned three positive cases
before stopping at the first ten-point shadow. The changed probe
records all eight outcomes; default replay performs certificate checks,
not a new classification or transporter search. Source transcription,
prior classifications and their retained subgroup-data premises are
acquired without additional census credit.

\begin{thebibliography}{9}
\bibitem{HS} D.~Leemans, The residually weakly primitive geometries of HS,
Australas. J. Combin. \textbf{33} (2005), 231--236, and the retained
23-page supplement dated 6 October 2003.
\bibitem{LPP} C.~Lef\`evre-Percsy, N.~Percsy and D.~Leemans,
New geometries for finite groups and polytopes,
Bull. Belg. Math. Soc. Simon Stevin \textbf{7} (2000), 583--610.
\end{thebibliography}
\end{document}
